\subsection{Automatic data-driven selection of representation and graph cutoff}
\label{sec:automatic_selection}

To avoid heuristic selection rules, the choice of semantic representation and sparsification cutoff
is treated as a joint empirical model-selection problem. The key principle is to separate
harmonization from graph construction. Step~2 remains a candidate-generation stage for possible
harmonization, whereas graph selection is performed only after the weighted graphs of Step~6
have been built. Until reviewed equivalence labels are available, the graph pipeline should
operate on a fixed no-merge baseline.

Let
\[
\mathcal{L}
=
\{\texttt{pooled},\texttt{track},\texttt{section},\texttt{programme},\texttt{programme\_year}\}
\]
denote the set of aggregation levels, and let
\[
\mathcal{R}
=
\{\texttt{name\_only},\texttt{concat}\}
\]
denote the set of candidate semantic representations. For each level
\(\ell \in \mathcal{L}\), let \(\mathcal{A}_{\ell}\) denote the set of analytical units at that
level. For each unit \(A \in \mathcal{A}_{\ell}\), let \(V_A\) denote the node universe fixed from
the Step~5 graph-ready node construction under the no-merge baseline. Both representations are
evaluated on this same frozen node set so that representation comparison is not confounded with
changes in the node universe.

\paragraph{Candidate graph families.}
For each level \(\ell\), unit \(A \in \mathcal{A}_{\ell}\), and representation
\(r \in \mathcal{R}\), Step~6 provides a weighted graph
\[
G^{(\ell,r)}_{A}
=
\bigl(V_A,E^{(\ell,r)}_A,w^{(\ell,r)}_A\bigr),
\]
where \(w^{(\ell,r)}_A(i,j)\) is the cosine similarity between nodes \(i\) and \(j\) under
representation \(r\). The candidate threshold set for level \(\ell\) and representation \(r\) is
defined as the set of distinct observed edge weights across all units at that level:
\[
\Theta_{\ell,r}
=
\left\{
w^{(\ell,r)}_A(i,j)
:
A \in \mathcal{A}_{\ell},\;
\{i,j\} \in E^{(\ell,r)}_A
\right\}.
\]
For any threshold \(\tau \in \Theta_{\ell,r}\), the corresponding sparse graph is
\[
H^{(\ell,r,\tau)}_A
=
\bigl(
V_A,\,
E^{(\ell,r,\tau)}_A,\,
w^{(\ell,r,\tau)}_A
\bigr),
\]
with
\[
E^{(\ell,r,\tau)}_A
=
\left\{
\{i,j\} \in E^{(\ell,r)}_A
:
w^{(\ell,r)}_A(i,j) \ge \tau
\right\},
\]
where surviving edges retain their original weights.

\paragraph{Cross-validated resampling at the occurrence level.}
Selection is based on out-of-sample performance, not on fixed topology rules. For each unit
\(A \in \mathcal{A}_{\ell}\), the source occurrences contributing to Step~5 are partitioned into
\(K\) folds. For fold \(k \in \{1,\dots,K\}\), let
\(\mathcal{O}^{\mathrm{train}}_{A,k}\) and \(\mathcal{O}^{\mathrm{test}}_{A,k}\) denote the
training and holdout occurrence sets. The node tables of Step~5 and the weighted graphs of
Step~6 are then rebuilt separately on the training and holdout occurrences.

To estimate the stability of edge retention under threshold \(\tau\), let
\(b \in \{1,\dots,B\}\) index bootstrap resamples drawn from
\(\mathcal{O}^{\mathrm{train}}_{A,k}\). For an unordered node pair
\(e=\{i,j\}\subseteq V_A\), define the training bootstrap edge-retention probability
\[
p^{(\ell)}_{e,A,k}(r,\tau)
=
\frac{1}{B}
\sum_{b=1}^{B}
\mathbf{1}
\!\left(
w^{(\ell,r,b)}_{e,A,k,\mathrm{train}} \ge \tau
\right),
\]
where \(w^{(\ell,r,b)}_{e,A,k,\mathrm{train}}\) is the bootstrap training edge weight for pair
\(e\). If a node is absent in a resample, the corresponding edge weight is treated as zero.

The holdout realization of the same thresholded edge is
\[
y^{(\ell)}_{e,A,k}(r,\tau)
=
\mathbf{1}
\!\left(
w^{(\ell,r)}_{e,A,k,\mathrm{test}} \ge \tau
\right).
\]

\paragraph{Proper scoring rule.}
For binary edge-retention outcomes, predictive performance is measured by the log-loss
\[
L(y,p)
=
-y\log p - (1-y)\log(1-p).
\]
Let
\[
\bar y^{(\ell)}_{A,k}(r,\tau)
=
\frac{1}{|P_A|}
\sum_{e \in P_A}
y^{(\ell)}_{e,A,k}(r,\tau),
\]
where \(P_A=\{\{i,j\}: i,j \in V_A,\; i<j\}\) is the set of all unordered node pairs in unit \(A\).
The quantity \(\bar y^{(\ell)}_{A,k}(r,\tau)\) is the holdout base-rate benchmark for that
unit-fold combination.

The cross-validated skill score for level \(\ell\), representation \(r\), and threshold \(\tau\)
is then
\[
\mathrm{Skill}_{\ell}(r,\tau)
=
1
-
\frac{
\sum\limits_{A \in \mathcal{A}_{\ell}}
\sum\limits_{k=1}^{K}
\sum\limits_{e \in P_A}
L\!\left(
y^{(\ell)}_{e,A,k}(r,\tau),
p^{(\ell)}_{e,A,k}(r,\tau)
\right)
}{
\sum\limits_{A \in \mathcal{A}_{\ell}}
\sum\limits_{k=1}^{K}
\sum\limits_{e \in P_A}
L\!\left(
y^{(\ell)}_{e,A,k}(r,\tau),
\bar y^{(\ell)}_{A,k}(r,\tau)
\right)
}.
\]
A larger value of \(\mathrm{Skill}_{\ell}(r,\tau)\) indicates better out-of-sample predictive
performance relative to the holdout edge-rate baseline.

\paragraph{Joint automatic selection.}
For each aggregation level \(\ell\), the selected representation and cutoff are chosen jointly as
\[
(r_{\ell}^{*},\tau_{\ell}^{*})
=
\arg\max_{r \in \mathcal{R},\;\tau \in \Theta_{\ell,r}}
\mathrm{Skill}_{\ell}(r,\tau).
\]
This yields one empirically selected graph-construction rule per level. The binary sparse graph
exported for level \(\ell\) is therefore based on the pair
\((r_{\ell}^{*},\tau_{\ell}^{*})\).

\paragraph{Decision rule for whether sparsification is supported.}
Thresholding should not be forced when the data do not support it. If
\[
\max_{r \in \mathcal{R},\;\tau \in \Theta_{\ell,r}}
\mathrm{Skill}_{\ell}(r,\tau)
\le 0,
\]
then no thresholded graph outperforms the trivial holdout base-rate model. In that case, the
automatic decision for level \(\ell\) is to retain the full weighted graph and not export a binary
sparse graph for that level.

\paragraph{Interpretation.}
This procedure eliminates the need for fixed admissibility rules based on largest-component share,
isolate share, susceptibility tie-breaking, or hand-set structural-versus-parsimony weights. The
representation and cutoff are selected directly from out-of-sample performance, and the decision
is made separately for each aggregation level rather than once in a pooled reference graph.

\paragraph{Workflow implications.}
Under this design, the workflow is modified as follows.
\begin{enumerate}
    \item Step~2 is used only for candidate generation for harmonization and no longer selects
    graph cutoffs.
    \item Steps~3 and~4 remain separate from graph construction. Until reviewed pair labels
    exist, the graph pipeline uses the no-merge baseline.
    \item Step~5 formally defines all graph levels:
    \(\texttt{pooled}\), \(\texttt{track}\), \(\texttt{section}\),
    \(\texttt{programme}\), and \(\texttt{programme\_year}\).
    \item Step~6 builds weighted graphs for both \texttt{name\_only} and \texttt{concat}
    on the same node universe.
    \item Steps~7 and~8 are replaced by one level-wise automatic model-selection step that
    jointly selects representation and cutoff.
    \item Step~9 uses the selected pair \((r_{\ell}^{*},\tau_{\ell}^{*})\) for each level.
\end{enumerate}

\paragraph{Recommended outputs.}
The automatic-selection step should write at least the following files:
\begin{itemize}
    \item \texttt{07\_levelwise\_model\_selection\_surface.csv}
    \item \texttt{07\_levelwise\_selected\_models.csv}
    \item \texttt{07\_edge\_selection\_probabilities\_<level>.csv}
\end{itemize}
The file \texttt{07\_levelwise\_model\_selection\_surface.csv} should contain one row per
\((\ell,r,\tau)\) combination, with columns such as
\texttt{level}, \texttt{representation}, \texttt{threshold},
\texttt{cv\_logloss}, \texttt{cv\_skill}, \texttt{n\_units},
\texttt{n\_candidate\_edges}, and \texttt{selected\_flag}. The file
\texttt{07\_levelwise\_selected\_models.csv} should contain one row per aggregation level with
the selected representation, selected threshold, cross-validated skill, and a logical field
indicating whether sparsification was supported. The file
\texttt{07\_edge\_selection\_probabilities\_<level>.csv} should report edge-level retention
probabilities under the selected model so that graph uncertainty remains visible in downstream
analysis.

\paragraph{Relation to later harmonization.}
Once reviewed equivalence labels are available from the harmonization workflow, an automatic
pair-classification model may be trained separately for synonym detection. However, that
harmonization model should remain distinct from the graph-selection procedure above. Pair
equivalence and graph sparsification are different inferential tasks and should not be collapsed
into a single cosine-cutoff rule.